\documentclass{situationist-preprint}

\title{Three Dimensions of Power Revisited}
\msid{POL-10-a}
\msfamily{POL-10}
\msdiscipline{Political Philosophy}
\mstype{Principal manuscript}
\mskeywords{power; Lukes; agenda-setting; preference formation; Dahl; domination}

\begin{document}
\maketitle

\begin{abstract}
What is power? The behavioral tradition identifies it with prevailing in observable conflicts over decisions; its critics add the capacity to keep issues off the agenda and the capacity to shape what people want. We argue that power is the capacity to shape agendas and preferences as well as outcomes. The argument has two steps. First, power is a capacity, not its exercise, and its object is which outcomes come about. Second, a principle of route-neutrality: if two capacities give an agent the same control over which outcomes come about, it cannot matter to whether either is power that one works through decisions and the other through the agenda or through preferences. Since control over agendas and preferences yields control over outcomes, it is power. We then depart from Lukes in one respect: power so understood does not require that it be exercised against anyone's real interests. The contrast between interest-respecting and interest-violating power belongs to the theory of domination, not to the concept of power, and separating them disarms the main objections to the third dimension.
\end{abstract}

\section{The problem}

\citet[ch.~1, §16]{Weber1922} defined power as the chance of an actor to carry out his own will within a social relationship despite resistance. \citet{Dahl1957} made the idea operational: $A$ has power over $B$ to the extent that $A$ can get $B$ to do something $B$ would not otherwise do. Dahl's followers studied power by observing who prevailed when decisions were contested.

\citet{Bachrach1962} objected that power has a second face. An actor who keeps an issue from ever being decided, by controlling what reaches the agenda or by mobilizing what \citet{Schattschneider1960} called the bias of the system, exercises power without prevailing in any observable contest. \citet{Crenson1971} found that air pollution did not become a political issue in a city dominated by a single steel firm, though it did in comparable cities; the firm's power showed in the absence of a decision. \citet{Lukes1974,Lukes2005} added a third dimension: $A$ may exercise power over $B$ by shaping $B$'s wants, so that no conflict arises at all. \citet{Gaventa1980} applied the three dimensions to an Appalachian valley in which miners acquiesced in conditions that harmed them.

The third dimension has been resisted on two grounds. It seems unobservable, since the counterfactual preferences $B$ would have had are not displayed; and Lukes defined it by reference to $B$'s real interests, which seemed to commit the analyst to knowing better than $B$ what $B$ wants. \citet{Morriss2002} argued, separately, that power is an ability, not an exercise of it, and Lukes accepted the point in the second edition \citep{Lukes2005}. We take Morriss's correction on board and show that, once it is, the three dimensions follow from a single principle, while the appeal to real interests can be dropped.

\section{Framework}

\begin{definition}[Power]\label{d:power}
An agent, individual or collective, has \emph{power} with respect to a range of outcomes to the extent that it has the capacity to make a difference to which of those outcomes come about.
\end{definition}

\begin{definition}[Routes]\label{d:routes}
An agent's capacity to make a difference to outcomes works through the \emph{decision route} if it consists in prevailing when alternatives are decided; through the \emph{agenda route} if it consists in determining which alternatives are put to decision; and through the \emph{preference route} if it consists in determining what those who decide want.
\end{definition}

\begin{principle}[Capacity]\label{p:capacity}
Power is a capacity, which may exist without being exercised.
\end{principle}

\begin{principle}[Route-neutrality]\label{p:route}
If two capacities give an agent the same degree of control over which outcomes in a range come about, then either both are power with respect to that range or neither is, whatever routes they work through.
\end{principle}

Principle~\ref{p:capacity} is Morriss's. The rich have power to buy votes whether or not they buy any, and an army has power to occupy a city whether or not it marches. Principle~\ref{p:route} expresses the point of the concept. We attribute power to agents in order to know whom to hold responsible for outcomes, whom to fear, and whom to bargain with \citep{Morriss2002}. For each of these purposes what matters is the degree of control over outcomes, not the mechanism. A legislature that would vote down any measure an industry opposes and a committee that ensures no such measure reaches the legislature leave the industry with the same control over regulation, and its opponents have the same reason to fear it.

\begin{assumption}\label{a:agenda}
For many ranges of outcomes, there are agents who, by determining which alternatives are put to decision, have as much control over which outcomes come about as they would have by prevailing in the decisions.
\end{assumption}

\begin{assumption}\label{a:pref}
For many ranges of outcomes, there are agents who, by determining what those who decide want, have as much control over which outcomes come about as they would have by prevailing in the decisions.
\end{assumption}

Assumption~\ref{a:agenda} is established by the literature on agenda control. Whoever sets the order of pairwise votes can often determine the result \citep{Riker1982}, and whoever determines that an issue is not raised secures the status quo as reliably as whoever defeats the proposal to change it \citep{Crenson1971}. Assumption~\ref{a:pref} is established by the ordinary workings of education, advertising and political communication. An agent who can bring a population to want a policy obtains that policy through the population's own decisions; the control is as real as if the agent had imposed it.

\section{The argument}

\begin{lemma}\label{l:decision}
The capacity to prevail in decisions is power.
\end{lemma}

\begin{proof}
To prevail in a decision between alternatives is to determine which of them comes about. By Definition~\ref{d:power} and Principle~\ref{p:capacity}, the capacity to do so is power with respect to the outcomes decided, whether or not it is exercised.
\end{proof}

\begin{theorem}\label{t:main}
Power includes the capacity to shape agendas and the capacity to shape preferences, as well as the capacity to determine outcomes by prevailing in decisions.
\end{theorem}

\begin{proof}
By Lemma~\ref{l:decision}, the capacity to prevail in decisions over a range of outcomes is power. By Assumption~\ref{a:agenda}, agents who control the agenda for many ranges have the same degree of control over outcomes as they would by prevailing; by Principle~\ref{p:route}, their capacity is power. By Assumption~\ref{a:pref} and Principle~\ref{p:route} likewise, the capacity to shape what deciders want is power.
\end{proof}

Where agenda control or preference shaping yields less control than prevailing would, Principle~\ref{p:route} does not apply directly, but Definition~\ref{d:power} does: the capacity still makes a difference to which outcomes come about, and it is power to the degree that it does. Theorem~\ref{t:main} is thus not a stipulation extending the word \enquote{power}; it follows from what the concept is for.

\begin{corollary}\label{c:interests}
Whether a capacity is power does not depend on whether its exercise would be contrary to the interests of those affected.
\end{corollary}

\begin{proof}
Definition~\ref{d:power} and Principles~\ref{p:capacity} and~\ref{p:route} make no reference to interests. A capacity to bring about outcomes that serve those affected gives as much control over outcomes as a capacity to bring about outcomes that harm them.
\end{proof}

Corollary~\ref{c:interests} departs from Lukes, for whom $A$ exercises power over $B$ when $A$ affects $B$ contrary to $B$'s interests \citep[ch.~4]{Lukes2005}. Lukes's definition marks a real distinction, but it is the distinction between power and domination: the capacity to interfere arbitrarily with another's choices, in the republican sense of \citet{Pettit1997}. A teacher who brings a pupil to want to learn has power over the pupil's preferences; whether the teacher dominates the pupil depends on whether the power is checked and answerable, which is a further question.

\section{Objections}

\begin{objection}[Unobservability]
Agenda and preference power leave no trace in observable conflict, so claims about them cannot be tested \citep{Polsby1963}.
\end{objection}

\begin{reply}
Every attribution of power, on Principle~\ref{p:capacity}, is a counterfactual claim, including attributions of decision power that has not been exercised. Counterfactuals are tested by comparison. Crenson compared cities that differed in the presence of a dominant polluter; Gaventa traced how quiescence in one valley gave way to challenge when the means of control weakened. These studies are as testable as the behavioral studies they criticize, and on the same principles.
\end{reply}

\begin{objection}[Structure, not agency]
Agendas and preferences are shaped by impersonal structures, by language, markets and institutions, rather than by agents. On a view such as Foucault's, power is everywhere and therefore no one's \citep{Foucault1975}.
\end{objection}

\begin{reply}
Much shaping of agendas and preferences is impersonal, and Definition~\ref{d:power} does not count it as anyone's power, since power is an agent's capacity. But much is not: the editor who decides what is reported, the firm that funds the research, the state that writes the curriculum, have capacities to shape agendas and preferences that they could exercise otherwise. Theorem~\ref{t:main} concerns these. That some influence on preferences is no one's power does not show that none is someone's.
\end{reply}

\begin{objection}[Over-extension]
On the present account parents, teachers, preachers and advertisers all have power, and the concept loses its critical edge.
\end{objection}

\begin{reply}
They do have power, and this is a merit of the account rather than a defect. The critical edge belongs to the theory of domination and of legitimacy, which asks whether a given power is checked, answerable and justifiable to those subject to it (POL-02). It is precisely because preference-shaping power is power that it falls within that theory's scope. The account also bears on POL-20, which requires equal access to public debate: unequal capacity to set the public agenda is unequal power, whatever the formal equality of votes, and it is relevant to whether democratic decisions have the intrinsic value POL-04 ascribes to them.
\end{reply}

\begin{objection}[Paternalism]
To say that $A$ shaped $B$'s preferences seems to imply that $B$'s preferences are not really $B$'s own, and to license others to override them.
\end{objection}

\begin{reply}
It implies neither. All preferences are shaped by someone or something, and a preference does not cease to be one's own because one acquired it from a teacher. By Corollary~\ref{c:interests}, the attribution of preference power is independent of any claim about $B$'s real interests, and so licenses nothing about overriding $B$'s preferences. Whether a shaped preference is autonomous is a question about the positive dimension of freedom (POL-09), not about power.
\end{reply}

\section{Conclusion}

Power is an agent's capacity to make a difference to which outcomes come about. Since the capacity to set agendas and the capacity to shape preferences can give an agent as much control over outcomes as the capacity to prevail in decisions, and since the route by which control is exercised does not bear on whether it is power, power includes all three. Lukes's three dimensions are therefore dimensions of power, but his tie between the third dimension and real interests is better placed in the theory of domination.

\begin{thebibliography}{99}\small
\bibitem[Bachrach and Baratz(1962)]{Bachrach1962} Bachrach, P. and Baratz, M.~S. (1962). Two faces of power. \emph{American Political Science Review}, 56(4), 947--952.
\bibitem[Crenson(1971)]{Crenson1971} Crenson, M.~A. (1971). \emph{The Un-Politics of Air Pollution: A Study of Non-Decisionmaking in the Cities}. Johns Hopkins Press.
\bibitem[Dahl(1957)]{Dahl1957} Dahl, R.~A. (1957). The concept of power. \emph{Behavioral Science}, 2(3), 201--215.
\bibitem[Foucault(1975)]{Foucault1975} Foucault, M. (1975). \emph{Surveiller et punir: Naissance de la prison}. Gallimard. Trans. A. Sheridan as \emph{Discipline and Punish}, Pantheon, 1977.
\bibitem[Gaventa(1980)]{Gaventa1980} Gaventa, J. (1980). \emph{Power and Powerlessness: Quiescence and Rebellion in an Appalachian Valley}. University of Illinois Press.
\bibitem[Lukes(1974)]{Lukes1974} Lukes, S. (1974). \emph{Power: A Radical View}. Macmillan.
\bibitem[Lukes(2005)]{Lukes2005} Lukes, S. (2005). \emph{Power: A Radical View}. 2nd edn. Palgrave Macmillan.
\bibitem[Morriss(2002)]{Morriss2002} Morriss, P. (2002). \emph{Power: A Philosophical Analysis}. 2nd edn. Manchester University Press.
\bibitem[Pettit(1997)]{Pettit1997} Pettit, P. (1997). \emph{Republicanism: A Theory of Freedom and Government}. Oxford University Press.
\bibitem[Polsby(1963)]{Polsby1963} Polsby, N.~W. (1963). \emph{Community Power and Political Theory}. Yale University Press.
\bibitem[Riker(1982)]{Riker1982} Riker, W.~H. (1982). \emph{Liberalism against Populism}. W.~H. Freeman.
\bibitem[Schattschneider(1960)]{Schattschneider1960} Schattschneider, E.~E. (1960). \emph{The Semisovereign People}. Holt, Rinehart and Winston.
\bibitem[Weber(1922)]{Weber1922} Weber, M. (1922). \emph{Wirtschaft und Gesellschaft}. Mohr. Trans. as \emph{Economy and Society}, ed. G. Roth and C. Wittich, University of California Press, 1978.
\end{thebibliography}

\end{document}
